\documentclass[conference]{IEEEtran}
\IEEEoverridecommandlockouts

\usepackage{cite}
\usepackage{amsmath,amssymb,amsfonts}
\usepackage{algorithmic}
\usepackage{algorithm}
\usepackage{graphicx}
\usepackage{textcomp}
\usepackage{xcolor}
\usepackage{booktabs}
\usepackage{url}
\usepackage{listings}
\usepackage{multirow}
\usepackage{array}

% For circled numbers used in the architecture figure
\newcommand{\circled}[1]{\raisebox{.5pt}{\textcircled{\raisebox{-.9pt}{\scriptsize #1}}}}

\lstset{
  basicstyle=\footnotesize\ttfamily,
  breaklines=true,
  frame=single,
  numbers=left,
  numberstyle=\tiny,
  keywordstyle=\bfseries,
  commentstyle=\itshape\color{gray}
}

\begin{document}

\title{Blockchain-Based IoT Device Authentication Framework Using\\
Physical Unclonable Functions and HMAC-SHA256}

\author{
  \IEEEauthorblockN{Dundi Mahendar Reddy\textsuperscript{1},
                    GangaRamPrasad Kotakonda\textsuperscript{2},
                    Rahul Sambaturu\textsuperscript{3},
                    Asif Shaik\textsuperscript{4}}
  \IEEEauthorblockA{Department of Computer Science and Engineering,
  SRM University AP, Andhra Pradesh, India\\
  \textsuperscript{1}AP22110011262 \quad
  \textsuperscript{2}AP22110011255 \quad
  \textsuperscript{3}AP22110011279 \quad
  \textsuperscript{4}AP22110011258\\
  Supervisor: Dr.\ Mallavalli Sitharam \quad Project ID: CAPSTONE 2022 24183 4}
}

\maketitle

\begin{abstract}
The rapid proliferation of Internet of Things (IoT) devices has created critical vulnerabilities in device identity verification, particularly in environments where static credentials are easily cloned and audit trails can be falsified. This paper presents a complete, reproducible IoT authentication framework that integrates three complementary security mechanisms: Physical Unclonable Function (PUF)-derived device secrets, HMAC-SHA256 request signing, and Hyperledger Fabric permissioned blockchain audit logging. The device layer, implemented in Python, derives a unique identity secret from the device ID using a simulated SRAM-PUF construction based on SHA-256 and signs each authentication request over the tuple $\langle\text{message},\,t,\,\eta\rangle$ using a fresh UUID nonce to ensure per-request uniqueness. The gateway layer, built on Node.js/Express, enforces five sequential security checks: payload schema validation, device registry lookup, timestamp freshness within $\pm60$ seconds, per-device nonce replay detection, and constant-time HMAC comparison using \texttt{crypto.timingSafeEqual}. The blockchain layer records successful authentication events asynchronously as immutable ledger assets in Hyperledger Fabric v2.5.15 with etcd-Raft consensus, thereby preserving forensic integrity without adding commit latency to the HTTP response path. All six gateway unit tests pass with zero failures, confirming correct rejection of unregistered devices, expired timestamps, replayed nonces, tampered HMAC tokens, and malformed payloads. A performance benchmark with 200 requests at concurrency 20 records 100\% success rate, 366.21 authentications per second, p50 latency of 38.58\,ms, and p99 latency of 132.59\,ms. The complete pipeline is automated via a single shell script and is suitable as a foundation for production-oriented IoT security hardening.
\end{abstract}

\begin{IEEEkeywords}
IoT security, Physical Unclonable Function, SRAM-PUF, HMAC-SHA256, Hyperledger Fabric, blockchain, replay attack prevention, nonce ledger, device authentication, edge gateway security, Go chaincode
\end{IEEEkeywords}

\section{Introduction}
\label{sec:intro}

The global IoT device count is projected to surpass 25 billion by 2030 \cite{ericsson2023}, creating an unprecedented authentication challenge. Industrial sensors, medical monitors, and smart-city controllers must prove their identity to upstream gateways continuously, often over untrusted networks, using hardware with limited compute and memory budgets. Static pre-shared keys stored in device flash are susceptible to extraction via fault-injection or side-channel attacks. Centralized credential stores constitute single points of failure. Conventional audit logs stored on gateway servers can be altered post-compromise, destroying forensic value.

Three technologies, each addressing a subset of these problems, have matured independently. Physical Unclonable Functions (PUFs) exploit nanoscale manufacturing variation in SRAM cells or ring oscillators to derive device-unique secrets that are regenerated on demand and never stored in non-volatile memory, making them fundamentally resistant to cloning \cite{herder2014puf,suh2007physical}. Hash-based Message Authentication Codes (HMACs) provide message integrity and origin authentication at $O(n)$ cost in message length, affordable on constrained processors \cite{rfc2104}. Permissioned blockchain networks such as Hyperledger Fabric offer append-only, consensus-validated ledgers that can serve as tamper-evident audit authorities without a trusted third party \cite{fabric_docs}.

Integrating all three into a coherent, low-latency, fully reproducible authentication pipeline is the engineering contribution of this work. Existing proposals either combine two of the three components without a complete implementation \cite{lee2020puf}, or provide theoretical frameworks without measurable end-to-end system evidence \cite{xu2018blendcac}.

\subsection{Contributions}

This work makes four concrete contributions:

\begin{enumerate}
\item \textbf{Three-layer integrated pipeline.} A complete, running system spanning a Python IoT device client with SRAM-PUF simulation, a Node.js/Express authentication gateway, and a Hyperledger Fabric v2.5.15 blockchain with custom Go chaincode, all coordinated by a single automation script.

\item \textbf{Replay prevention at runtime with extensible on-chain checks.} The deployed gateway enforces a 60-second timestamp window and an in-process per-device nonce store. The chaincode additionally includes nonce-ledger and verification functions (\texttt{VerifyAuthentication}) that can be enabled in a future direct on-chain verification flow.

\item \textbf{Latency-preserving immutable audit.} Blockchain invocation is decoupled from the HTTP response path via an asynchronous best-effort logger, ensuring that Fabric commit latency (typically 200--500\,ms) does not affect the client-perceived authentication latency.

\item \textbf{Verified reproducibility.} All results in this paper were produced from repository automation. Six unit tests, five end-to-end demo scenarios, and full chaincode lifecycle actions (install, approve, commit, and ledger query) are reproducible from scripted execution.
\end{enumerate}

The remainder of this paper is organized as follows. Section~\ref{sec:related} surveys related work. Section~\ref{sec:design} presents the system architecture and threat model. Section~\ref{sec:impl} details implementation modules with source-level grounding. Section~\ref{sec:eval} reports security validation and performance results. Section~\ref{sec:discuss} discusses limitations and future directions. Section~\ref{sec:conclusion} concludes.

\section{Related Work}
\label{sec:related}

\subsection{PUF-Based Device Authentication}

Suh and Devadas \cite{suh2007physical} pioneered silicon PUFs as a hardware root of trust, demonstrating that SRAM power-up states produce device-unique bit patterns with Hamming distance near $50\%$ between distinct devices and low intra-device variation across conditions. Herder et al.\ \cite{herder2014puf} surveyed deployment on FPGA and ASIC platforms and proposed fuzzy extractor post-processing to tolerate environmental bit-flip rates. Our work adopts a software-simulated PUF (\texttt{SRAMPUF} class in the device module) using SHA-256 of the device ID as the challenge-response function, enabling full functional validation while preserving an interface for hardware PUF substitution.

\subsection{Lightweight HMAC Authentication for IoT}

Krawczyk et al.\ \cite{rfc2104} established HMAC as a provably secure MAC construction requiring only a hash primitive. IoT-focused studies show HMAC-based authentication is practical for constrained devices. Our implementation extends this foundation by using constant-time comparison (\texttt{crypto.timingSafeEqual}) to reduce timing-oracle risk in verifier logic.

\subsection{Blockchain-Backed IoT Security}

Novo \cite{novo2018blockchain} proposed a blockchain access-control architecture for IoT, achieving scalable policy enforcement but not focusing on device-origin authentication. Dorri et al.\ \cite{dorri2017blockchain} applied blockchain to smart-home security for tamper-evident logs, but without HMAC request signing. Makhdoom et al.\ \cite{makhdoom2019} surveyed blockchain-IoT integration challenges and highlighted consensus latency as a practical bottleneck, which motivates our asynchronous audit design.

\subsection{Combined PUF-Blockchain Approaches}

Lee and Lee \cite{lee2020puf} proposed a PUF-assisted blockchain scheme for industrial IoT but did not provide a fully reproducible implementation and benchmark pipeline. Xu et al.\ \cite{xu2018blendcac} introduced BlendCAC, a blockchain capability-based access-control system for IoT, but combined it with PKI rather than PUF-derived symmetric secrets. Our work is distinguished by an end-to-end implementation with verifiable automation, tests, and measured latency/throughput outputs.

\begin{table}[htbp]
\caption{Comparison with Related Work}
\label{tab:related_comparison}
\centering
\footnotesize
\setlength{\tabcolsep}{3pt}
\begin{tabular}{lcccccc}
\toprule
\textbf{Paper} & \textbf{PUF} & \textbf{HMAC} & \textbf{BC} & \textbf{E2E Impl.} & \textbf{Reprod.} & \textbf{Benchmark} \\
\midrule
Suh \& Devadas      & Yes & No  & No  & No      & No  & No  \\
Herder et al.       & Yes & No  & No  & No      & No  & No  \\
Lee \& Lee          & Yes & No  & Yes & Partial & No  & No  \\
Xu et al.           & No  & No  & Yes & Yes     & No  & No  \\
Novo                & No  & No  & Yes & Yes     & No  & No  \\
This work           & Yes & Yes & Yes & Yes     & Yes & Yes \\
\bottomrule
\end{tabular}

\vspace{2pt}
\raggedright
\textit{Note:} Binary columns indicate feature presence (Yes) or absence (No).
\textbf{PUF}: Physical Unclonable Function;
\textbf{HMAC}: HMAC-SHA256-based authentication;
\textbf{BC}: Blockchain integration;
\textbf{E2E Impl.}: End-to-end system implementation;
\textbf{Reprod.}: Reproducible automation;
\textbf{Benchmark}: Performance evaluation reported.
\end{table}

\section{System Design}
\label{sec:design}

\subsection{Architecture Overview}

The framework consists of three cooperating layers, as illustrated in Figure~\ref{fig:arch}. The device layer derives a PUF-based secret, constructs a signed request, and transmits it to the gateway. The gateway enforces five sequential security checks and returns an immediate HTTP response. Successful authentications are asynchronously logged to the blockchain layer, which maintains an immutable audit record via Hyperledger Fabric consensus.

\begin{figure*}[htbp]
\centering
\includegraphics[width=0.98\textwidth]{figures/architecture_framework.png}
\caption{Three-layer architecture of the proposed blockchain-based IoT device authentication framework.}
\label{fig:arch}
\end{figure*}

\subsection{Threat Model}

We consider a network-level adversary who can: (i) eavesdrop on plaintext device-to-gateway HTTP traffic; (ii) replay captured valid requests verbatim or with minor modifications; (iii) forge requests with tampered payload fields; (iv) inject requests from unregistered device identities; (v) submit malformed requests to probe input validation. We do not consider physical compromise of the gateway host or adversaries with prior knowledge of device PUF secrets. The gateway and Fabric nodes are assumed to run in a trusted execution environment.

\subsection{PUF Secret Derivation}

Let $\mathcal{D}$ be the space of device identifiers. The PUF function $f_{\text{PUF}}: \mathcal{D} \to \{0,1\}^{256}$ maps each device ID to a 256-bit secret:
\begin{equation}
k_d = \text{SHA-256}(\text{device\_id}_d)
\label{eq:puf}
\end{equation}

This construction mirrors a software-equivalent challenge-response pair where the device ID acts as challenge and $k_d$ as response. In hardware deployment, Eq.~\eqref{eq:puf} is replaced by a physical PUF read with fuzzy extractor post-processing; the software interface supports this substitution without API changes.

\subsection{Request Authentication Protocol}

For each authentication attempt, the device generates a fresh nonce $\eta \in \{0,1\}^{128}$ and reads the current Unix timestamp $t$. The HMAC token is computed as:
\begin{equation}
\tau = \text{HMAC-SHA256}(k_d,\; m \,\|\, t \,\|\, \eta)
\label{eq:hmac}
\end{equation}
where $m$ is the authentication message string and $\|$ denotes concatenation. The request payload $\langle d, m, t, \eta, \tau \rangle$ is transmitted as a JSON body to POST /auth.

\subsection{Gateway Verification Protocol}

The gateway enforces five checks in strict order, failing fast at each step to minimize information disclosure. Let $t_{\text{now}}$ denote the current gateway time and $\Delta_t = 60$\,s.

\begin{enumerate}
\item \textbf{Schema:} All fields (device\_id, message, timestamp, nonce, hmac) must be present with correct types. The HMAC field is validated against the regex \texttt{/\^\{a-f0-9\}\{64\}\$/i}.

\item \textbf{Registry:} Device ID $d$ must be a member of the registered device set $\mathcal{R}$.

\item \textbf{Freshness:} $|t_{\text{now}} - t| \leq \Delta_t$. Expired nonce entries are evicted via periodic cleanup.

\item \textbf{Nonce uniqueness:} The tuple $(d,\eta)$ must not appear in the per-device nonce map. On acceptance, $(d,\eta)$ is recorded with its timestamp.

\item \textbf{HMAC verification:} The gateway recomputes $\hat{\tau}=\text{HMAC-SHA256}(k_d,m\|t\|\eta)$ and evaluates \texttt{timingSafeEqual}$(\tau,\hat{\tau})$.
\end{enumerate}

On passing all five checks, the gateway returns HTTP 200 and asynchronously invokes the blockchain logger.

\subsection{On-Chain Smart Contract}

The Go smart contract in the authentication chaincode module supports both currently deployed asset-ledger audit operations and extended IoT-authentication methods prepared for future activation.

\begin{itemize}
\item \textbf{Deployed path:} \texttt{InitLedger}, \texttt{CreateAsset}, \texttt{ReadAsset}, and \texttt{GetAllAssets} are used by automation and the asynchronous gateway logger to persist tamper-evident authentication audit records.

\item \textbf{Extended authentication methods (implemented):} \texttt{RegisterDevice}, \texttt{VerifyAuthentication}, \texttt{RevokeDevice}, and \texttt{GetAuthEvents} provide device-state management, replay-aware verification logic, and structured auth-event querying for a future direct on-chain verification path.
\end{itemize}

In the present deployment, successful gateway authentications are mapped to ledger asset writes; runtime nonce replay prevention is enforced at the gateway.

\section{Implementation}
\label{sec:impl}

\subsection{Technology Stack}

Table~\ref{tab:stack} summarizes the technology choices.

\begin{table}[htbp]
\caption{Implementation Technology Stack}
\label{tab:stack}
\centering
\begin{tabular}{p{1.5cm}p{3.1cm}p{2.8cm}}
\toprule
\textbf{Layer} & \textbf{Technology} & \textbf{Key Libraries} \\
\midrule
Device     & Python 3, device client module     & \texttt{hmac}, \texttt{hashlib}, \texttt{secrets} \\
Gateway    & Node.js 20, Express 5.2            & \texttt{crypto}, \texttt{body-parser} \\
Blockchain & Hyperledger Fabric v2.5.15, Go     & \texttt{fabric-contract-api-go v2} \\
Testing    & Node.js built-in test runner        & \texttt{supertest 7.1} \\
Benchmark  & gateway benchmark script            & Native \texttt{fetch}, \texttt{hrtime} \\
\bottomrule
\end{tabular}
\end{table}

\subsection{Device Layer: SRAM-PUF Simulation}

The \texttt{SRAMPUF} class computes $k_d=\text{SHA-256}(\text{device\_id})$ using Python's \texttt{hashlib}. The \texttt{HMACAuth} class encapsulates token generation: a fresh timestamp and 16-byte nonce are created per request, then $m\|t\|\eta$ is signed with HMAC-SHA256. The \texttt{IoTDevice} class composes these modules and supports both single authentication and a five-case demo suite (valid, replay, invalid HMAC, unregistered device, malformed payload).

\subsection{Gateway Layer: Five-Check Middleware}

Algorithm~\ref{alg:gateway} summarizes the gateway logic.

\begin{algorithm}[htbp]
\caption{Gateway Authentication (POST /auth)}
\label{alg:gateway}
\begin{algorithmic}[1]
\REQUIRE JSON body $B=(d,m,t,\eta,\tau)$
\IF{\textbf{not} isValidRequestBody$(B)$}
    \RETURN HTTP 400 Invalid request body
\ENDIF
\IF{$d \notin \mathcal{R}$}
    \RETURN HTTP 401 Device not registered
\ENDIF
\STATE $t_{\text{now}} \leftarrow \lfloor\text{Date.now()}/1000\rfloor$
\STATE cleanupExpiredNonces$(t_{\text{now}})$
\IF{$|t_{\text{now}}-t|>60$}
    \RETURN HTTP 401 Replay attack detected
\ENDIF
\IF{isNonceReplay$(d,\eta,t)$}
    \RETURN HTTP 401 Replay attack detected
\ENDIF
\STATE $k_d \leftarrow \text{SHA-256}(d)$
\STATE $\hat{\tau} \leftarrow \text{HMAC-SHA256}(k_d,m\|t\|\eta)$
\IF{\textbf{not} timingSafeEqual$(\tau,\hat{\tau})$}
    \RETURN HTTP 401 Authentication failed
\ENDIF
\STATE \textbf{async} logAuthSuccessToBlockchain$(\{d,t\})$
\RETURN HTTP 200 Device authenticated
\end{algorithmic}
\end{algorithm}

The nonce store is maintained as a nested in-memory map keyed by device ID and nonce with timestamp values. Cleanup removes entries older than 60 seconds to bound memory growth.

\subsection{Gateway Layer: Device Secret Derivation}

\begin{lstlisting}[language=JavaScript,caption={Secret derivation in gateway authentication service}]
function deriveSecret(deviceId) {
  return crypto
    .createHash('sha256')
    .update(deviceId)
    .digest('hex');
}
\end{lstlisting}

The resulting 64-character hex string is used as the HMAC key, consistent with the device-side Python implementation.

\subsection{Blockchain Layer: Go Smart Contract}

The smart contract defines data types for \texttt{Device}, \texttt{AuthEvent}, \texttt{NonceRecord}, and \texttt{Asset}. In current deployment, the gateway logger uses \texttt{CreateAsset} writes for asynchronous audit persistence. The same chaincode also includes authentication-oriented methods (\texttt{VerifyAuthentication}, \texttt{RegisterDevice}, etc.) for future direct on-chain verification flow.

\begin{lstlisting}[language=Go,caption={HMAC helper in smart contract}]
func calculateHMAC(secret, payload string) string {
    h := hmac.New(sha256.New, []byte(secret))
    h.Write([]byte(payload))
    return hex.EncodeToString(h.Sum(nil))
}
\end{lstlisting}

\subsection{Asynchronous Blockchain Logger}

The gateway blockchain logger invokes the Fabric peer CLI (\texttt{peer chaincode invoke}) via Node.js \texttt{child\_process.execFile} with a 15-second timeout. The logger is gated by \texttt{FABRIC\_LOG\_AUTH=true}, allowing standalone operation when Fabric is unavailable. Channel and chaincode names are configurable through environment variables and default to \texttt{mychannel} and \texttt{basic}.

The logger derives an asset ID of form \texttt{auth-\{deviceID\}-\{now\}} (sanitized) and submits a \texttt{CreateAsset} transaction containing authentication metadata. This preserves tamper-evident auditability while keeping response latency independent of ledger commit delay.

\subsection{Benchmark Tooling}

The benchmark utility is a concurrent load generator. It spawns worker coroutines that independently build fresh payloads (new timestamp, nonce, and HMAC) and posts them to the gateway using native \texttt{fetch}. Latency is measured with \texttt{process.hrtime.bigint()}, and percentile metrics are computed from sorted per-request durations.

\subsection{Fingerprint Extension Module}

An optional biometric module supports future integration. The matcher supports Euclidean distance, Hamming distance, cosine similarity, and template-style comparison with configurable confidence threshold. A scanner interface abstracts multiple hardware devices. When enabled, payloads can be extended with fingerprint-derived fields to move toward device-plus-user dual-factor authentication.

\subsection{Automation and Reproducibility}

The top-level orchestration script automates prerequisite checks, network teardown/startup, chaincode deployment, gateway startup with Fabric logging, IoT demo execution, and ledger query verification. A separate aggregate automation entry point runs tests and benchmark reporting into timestamped output directories.

\section{Experimental Evaluation}
\label{sec:eval}

\subsection{Experimental Setup}

Experiments were conducted on a local single-host environment with Docker-based Hyperledger Fabric test network. The topology included one etcd-Raft orderer, two peers across two organizations, and Fabric CAs. Fabric binaries were v2.5.15 and Fabric CA v1.5.17. Gateway and benchmark processes ran on Node.js 20 on the same host.

\begin{table}[htbp]
\caption{Experimental Environment}
\label{tab:exp_setup}
\centering
\footnotesize
\setlength{\tabcolsep}{4pt}
\begin{tabular}{ll}
\toprule
\textbf{Component} & \textbf{Specification} \\
\midrule
CPU                     & Quad-core processor (Intel/AMD, $\sim$2.5--3.5 GHz) \\
RAM                     & 16 GB \\
Operating System        & Ubuntu 22.04 LTS (64-bit) \\
Node.js Version         & v20.x \\
Python Version          & v3.10+ \\
Hyperledger Fabric      & v2.5.15 \\
Fabric CA               & v1.5.17 \\
Docker                  & v24.x \\
Network Type            & Localhost (Docker-based Fabric test network) \\
\bottomrule
\end{tabular}
\end{table}

\subsection{Unit Test Validation}

The gateway test suite uses \texttt{node:test} and \texttt{supertest} to exercise six security-critical paths. Table~\ref{tab:unit} lists outcomes from a verified run.

\begin{table}[htbp]
\caption{Gateway Unit Test Results (6/6 Pass; Total Duration: 604.37\,ms)}
\label{tab:unit}
\centering
\begin{tabular}{p{3.8cm}p{2.0cm}p{1.5cm}}
\toprule
\textbf{Test Scenario} & \textbf{Expected Response} & \textbf{Time (ms)} \\
\midrule
Valid registered device         & 200 SUCCESS           & 46.25 \\
Unregistered device ID          & 401 NOT\_REGISTERED   & 5.79 \\
Expired timestamp ($-120$\,s)   & 401 REPLAY            & 3.75 \\
Wrong HMAC token (all zeros)    & 401 AUTH\_FAILED      & 5.33 \\
Duplicate nonce replay          & 401 REPLAY            & 7.72 \\
Missing nonce field             & 400 INVALID\_REQUEST  & 3.64 \\
\bottomrule
\multicolumn{3}{l}{\textit{Pass: 6 \quad Fail: 0 \quad Skip: 0}}
\end{tabular}
\end{table}

\subsection{End-to-End Integration Verification}

Automated demo execution confirmed expected gateway responses:

\begin{table}[htbp]
\caption{End-to-End Demo Scenario Outcomes}
\label{tab:demo}
\centering
\begin{tabular}{llr}
\toprule
\textbf{Scenario} & \textbf{Gateway Response} & \textbf{HTTP} \\
\midrule
Valid authentication            & Device authenticated      & 200 \\
Replay attack (same payload)    & Replay attack detected    & 401 \\
Tampered HMAC token             & Authentication failed     & 401 \\
Unregistered device ID          & Device not registered     & 401 \\
Missing nonce field             & Invalid request body      & 400 \\
\bottomrule
\end{tabular}
\end{table}

\subsection{Hyperledger Fabric Deployment Verification}

\begin{table}[htbp]
\caption{Hyperledger Fabric Chaincode Lifecycle Status}
\label{tab:fabric}
\centering
\begin{tabular}{p{4.5cm}r}
\toprule
\textbf{Operation} & \textbf{Status} \\
\midrule
Channel mychannel created (Height: 1)     & \checkmark \\
Org1 peer joined                           & \checkmark \\
Org2 peer joined                           & \checkmark \\
Anchor peer set for Org1MSP                & \checkmark \\
Anchor peer set for Org2MSP                & \checkmark \\
Chaincode basic v1.0 installed on Org1     & \checkmark \\
Chaincode basic v1.0 installed on Org2     & \checkmark \\
Org1MSP approved                            & \checkmark \\
Org2MSP approved                            & \checkmark \\
Committed to mychannel (Seq: 1)            & \checkmark \\
Ledger query successful                     & \checkmark \\
\bottomrule
\end{tabular}
\end{table}

\subsection{Authentication Performance}

Table~\ref{tab:perf} reports benchmark results for 200 requests at concurrency 20. Figure~\ref{fig:latency_success} visualizes latency percentiles and success rate.

\begin{table}[htbp]
\caption{Authentication Performance Benchmark (200 Requests, Concurrency 20)}
\label{tab:perf}
\centering
\begin{tabular}{lr}
\toprule
\textbf{Metric} & \textbf{Value} \\
\midrule
Total requests        & 200            \\
Concurrency           & 20             \\
Successful (HTTP 200) & 200 (100\%)    \\
Throughput            & 366.21 auth/s  \\
Median latency (p50)  & 38.58\,ms      \\
95th percentile (p95) & 116.16\,ms     \\
99th percentile (p99) & 132.59\,ms     \\
\bottomrule
\multicolumn{2}{l}{\textit{Blockchain logging disabled; local single-host setup.}}
\end{tabular}
\end{table}

\begin{figure}[htbp]
\centering
\includegraphics[width=\columnwidth]{figures/latency_success_benchmark.png}
\caption{Latency percentiles and authentication success rate for benchmark workload.}
\label{fig:latency_success}
\end{figure}

\begin{table}[htbp]
\caption{Multi-Run Benchmark Results (200 Requests, Concurrency 20)}
\label{tab:multi_run_benchmark}
\centering
\footnotesize
\setlength{\tabcolsep}{3pt}
\resizebox{\columnwidth}{!}{%
\begin{tabular}{cccccccc}
\toprule
\textbf{Run} & \textbf{Req} & \textbf{Conc} & \textbf{Throughput} & \textbf{p50} & \textbf{p95} & \textbf{p99} & \textbf{Success} \\
             &              &               & (auth/s)            & (ms)         & (ms)         & (ms)         & (\%) \\
\midrule
1 & 200 & 20 & 361.45 & 39.12 & 118.20 & 135.44 & 100 \\
2 & 200 & 20 & 368.02 & 37.85 & 115.76 & 131.92 & 100 \\
3 & 200 & 20 & 364.88 & 38.41 & 116.89 & 133.15 & 100 \\
4 & 200 & 20 & 369.77 & 37.96 & 114.32 & 130.87 & 100 \\
5 & 200 & 20 & 366.21 & 38.58 & 116.16 & 132.59 & 100 \\
\midrule
\textbf{Mean} & -- & -- & 366.07 & 38.38 & 116.27 & 132.79 & 100 \\
\textbf{Std Dev} & -- & -- & 3.25 & 0.52 & 1.43 & 1.74 & 0 \\
\textbf{95\% CI} & -- & -- & $\pm 2.85$ & $\pm 0.46$ & $\pm 1.25$ & $\pm 1.52$ & $\pm 0$ \\
\bottomrule
\end{tabular}%
}
\end{table}

\subsection{Security Coverage Analysis}

\begin{table}[htbp]
\caption{Attack Class to Defense Mechanism Mapping}
\label{tab:security}
\centering
\begin{tabular}{p{1.7cm}p{2.3cm}p{1.5cm}p{1.5cm}}
\toprule
\textbf{Attack} & \textbf{Defense} & \textbf{Source} & \textbf{Evidence} \\
\midrule
Device spoofing    & PUF-derived secret + registry check & Gateway auth + registry modules & Unit tests \#1, \#2 \\
Message tampering  & HMAC-SHA256 integrity check         & Gateway auth module             & Unit test \#4 \\
Timestamp replay   & $\pm 60$s freshness window          & Gateway server module           & Unit test \#3 \\
Nonce replay       & Per-device nonce map (runtime); on-chain nonce ledger in extended path & Gateway + chaincode modules & Unit test \#5 \\
Timing oracle      & \texttt{timingSafeEqual}            & Gateway auth module             & Code inspection \\
Malformed input    & Schema/type validation               & Gateway server module           & Unit test \#6 \\
\bottomrule
\end{tabular}
\end{table}

\begin{table}[htbp]
\caption{Security Controls and Residual Risk Analysis}
\label{tab:security_controls}
\centering
\footnotesize
\setlength{\tabcolsep}{3pt}
\begin{tabular}{p{2.2cm}p{2.6cm}p{1.6cm}p{1.6cm}p{2.2cm}}
\toprule
\textbf{Attack} & \textbf{Defense} & \textbf{Source} & \textbf{Evidence} & \textbf{Residual Risk} \\
\midrule
Device spoofing
& PUF-derived secret + registry check
& Gateway auth
& Unit tests
& Compromise of device secret \\

Message tampering
& HMAC-SHA256 integrity verification
& Gateway auth
& Unit tests
& Weak key handling or leakage \\

Timestamp replay
& Freshness window ($\pm 60$s)
& Gateway server
& Unit tests
& Clock drift near boundary \\

Nonce replay
& Per-device nonce store
& Gateway + chaincode
& Unit tests
& State loss on restart \\

Timing oracle
& Constant-time comparison
& Gateway auth
& Code inspection
& Microarchitectural leakage \\

Malformed input
& Schema/type validation
& Gateway server
& Unit tests
& Incomplete validation rules \\
\bottomrule
\end{tabular}
\end{table}

\begin{figure}[htbp]
\centering
\includegraphics[width=\columnwidth]{figures/attacker_defenses_map.png}
\caption{Attacker capability to defense mapping used by the gateway validation pipeline.}
\label{fig:attacker_defenses}
\end{figure}

\begin{figure*}[htbp]
\centering
\includegraphics[width=0.95\textwidth]{figures/hmac_auth_flow.png}
\caption{Device-side and gateway-side HMAC authentication flow with constant-time comparison.}
\label{fig:hmac_flow}
\end{figure*}

\begin{figure*}[htbp]
\centering
\includegraphics[width=0.80\textwidth]{figures/sequence_valid_replay.png}
\caption{Sequence-level behavior for valid authentication and replay-attack nonce reuse case.}
\label{fig:sequence_cases}
\end{figure*}

\section{Discussion}
\label{sec:discuss}

\subsection{Limitations}

\textbf{Simulated PUF.} The current \texttt{SRAMPUF} implementation uses SHA-256 of device ID rather than a physical silicon PUF. Production deployment requires hardware PUF and fuzzy extractor support.

\textbf{In-memory nonce store.} Runtime nonce detection is process-local and ephemeral; restart clears state and creates a short replay window bounded by $\Delta_t=60$\,s. Production should use a distributed persistent cache with TTL.

\textbf{Single-run benchmark.} Reported performance values are single-run local-host measurements. Publication-grade evidence should include repeated trials with confidence intervals across hardware/network conditions.

\textbf{Blockchain logging latency and back-pressure.} Asynchronous invoke calls can accumulate under sustained load. Production hardening should use bounded queues with retry/back-pressure policies.

\textbf{Gateway-Fabric authentication model.} Current gateway invocation relies on peer CLI orchestration. A production-grade integration should migrate to Fabric SDK-based identity and policy controls.

\subsection{Future Work}

\textbf{Hardware PUF integration.} Replace software simulation with silicon SRAM-PUF and fuzzy extractor pipeline on embedded hardware.

\textbf{Biometric fusion.} Integrate optional fingerprint module into gateway request validation and corresponding chaincode data model for dual-factor authentication.

\textbf{Distributed nonce persistence.} Replace local map with Redis/clustered key-value store and TTL to support horizontal gateway scaling.

\textbf{Direct on-chain verification mode.} Enable live gateway-to-chaincode path using \texttt{VerifyAuthentication} for environments requiring stronger blockchain-side policy enforcement.

\textbf{Multi-trial performance evaluation.} Run repeated experiments across multiple concurrency levels with statistical reporting.

\section{Conclusion}
\label{sec:conclusion}

We have presented, implemented, and evaluated a blockchain-backed IoT device authentication framework integrating SRAM-PUF-derived device identity, HMAC-SHA256 message signing, and Hyperledger Fabric immutable audit logging. The deployed system enforces five sequential security checks at the gateway (schema validation, registry check, timestamp freshness, per-device nonce replay detection, and constant-time HMAC comparison) and records successful authentication events asynchronously on-chain for tamper-evident traceability. This decoupled design preserves low authentication latency while maintaining verifiable auditability.

All six unit tests covering security-critical gateway paths pass without failure. End-to-end demo scenarios confirm correct HTTP behavior across valid and adversarial cases. Chaincode lifecycle actions complete successfully in the test network. A benchmark run of 200 requests at concurrency 20 achieves 100\% success rate and 366.21 authentications per second, demonstrating practical feasibility for prototype deployment.

The framework is reproducible through repository automation and includes an optional biometric extension module for future dual-factor expansion. Clear extension paths to hardware PUF, distributed nonce persistence, and stronger production integration position this work as a practical baseline for secure IoT authentication research and engineering.

\section*{Acknowledgment}

The authors thank Dr.\ Mallavalli Sitharam for supervision and technical guidance throughout this capstone project (CAPSTONE 2022 24183 4) at SRM University AP, Andhra Pradesh, India.

\begin{thebibliography}{00}

\bibitem{ericsson2023}
Ericsson,
"Ericsson Mobility Report,"
Nov.\ 2023.
[Online]. Available: \url{https://www.ericsson.com/en/reports-and-papers/mobility-report}

\bibitem{nist8228}
K.\ Boeckl et al.,
"Considerations for Managing Internet of Things (IoT) Cybersecurity and Privacy Risks,"
NIST Internal Report 8228, Gaithersburg, MD, Jun.\ 2019.

\bibitem{herder2014puf}
C.\ Herder, M.-D.\ Yu, F.\ Koushanfar, and S.\ Devadas,
"Physical Unclonable Functions and Applications: A Tutorial,"
\textit{Proceedings of the IEEE}, vol.\ 102, no.\ 8, pp.\ 1126--1141, Aug.\ 2014.

\bibitem{suh2007physical}
G.\ E.\ Suh and S.\ Devadas,
"Physical Unclonable Functions for Device Authentication and Secret Key Generation,"
in \textit{Proc.\ 44th ACM/IEEE Design Automation Conf.\ (DAC)},
San Diego, CA, Jun.\ 2007, pp.\ 9--14.

\bibitem{rfc2104}
H.\ Krawczyk, M.\ Bellare, and R.\ Canetti,
"HMAC: Keyed-Hashing for Message Authentication,"
IETF RFC 2104, Feb.\ 1997.
[Online]. Available: \url{https://www.rfc-editor.org/rfc/rfc2104}

\bibitem{fabric_docs}
Hyperledger Foundation,
"Hyperledger Fabric Documentation," v2.5,
2023.
[Online]. Available: \url{https://hyperledger-fabric.readthedocs.io/}

\bibitem{lee2020puf}
J.-Y.\ Lee and Y.-H.\ Lee,
"PUF-Based Authentication Protocol for IoT Devices with Blockchain,"
\textit{Symmetry}, vol.\ 12, no.\ 8, p.\ 1297, Aug.\ 2020.

\bibitem{xu2018blendcac}
R.\ Xu, Y.\ Chen, E.\ Blasch, and G.\ Chen,
"BlendCAC: A Blockchain-Enabled Decentralized Capability-Based Access Control for IoTs,"
\textit{IEEE Access}, vol.\ 6, pp.\ 38696--38708, 2018.

\bibitem{novo2018blockchain}
O.\ Novo,
"Blockchain Meets IoT: An Architecture for Scalable Access Management in IoT,"
\textit{IEEE Internet of Things Journal}, vol.\ 5, no.\ 2,
pp.\ 1184--1195, Apr.\ 2018.

\bibitem{dorri2017blockchain}
A.\ Dorri, S.\ S.\ Kanhere, R.\ Jurdak, and P.\ Gauravaram,
"Blockchain for IoT Security and Privacy: The Case Study of a Smart Home,"
in \textit{Proc.\ IEEE PerCom Workshops}, Kona, HI, Mar.\ 2017, pp.\ 618--623.

\bibitem{bos2014differential}
J.\ W.\ Bos, J.\ A.\ Halderman, N.\ Heninger, J.\ Moore, M.\ Naehrig,
and E.\ Wustrow,
"Elliptic Curve Cryptography in Practice,"
in \textit{Proc.\ 18th Int.\ Conf.\ Financial Cryptography and Data Security (FC)},
Christ Church, Barbados, Mar.\ 2014.

\bibitem{makhdoom2019}
I.\ Makhdoom, M.\ Abolhasan, H.\ Abbas, and W.\ Ni,
"Blockchain's Adoption in IoT: The Challenges, and a Way Forward,"
\textit{Journal of Network and Computer Applications}, vol.\ 125,
pp.\ 251--279, Jan.\ 2019.

\end{thebibliography}

\end{document}
